% =============================================================================
\chapter{Discussion}
\label{ch:discussion}
% =============================================================================

\section{Interpretation of the Results}
\label{sec:disc:interpret}

The results in \cref{ch:results} confirm the central hypothesis of this
work: a hands-and-pose landmark representation combined with a deep
temporal model is sufficient to discriminate hundreds of word-level signs
in both \gls{asl} and \gls{arsl} without resorting to heavy 3-D CNNs. The
biggest individual gains come from (a) extending the sequence window
from 30 to 48 frames and (b) augmenting hands with the 33-landmark pose,
both of which carry strong body-anchor information for signs such as
\emph{house}, \emph{I}, or \emph{you}.

\section{Strengths}
\label{sec:disc:strengths}

\begin{itemize}
  \item The unified pipeline minimises the engineering surface --- one
        feature extractor feeds two language pairs and two granularity
        levels.
  \item The custom \texttt{TemporalAttention} layer is trivially small
        ($<\,1$\,\% parameters) yet contributes measurable Top-5 gains
        (\cref{tab:ablate}).
  \item The system is genuinely deployed --- not a notebook demo --- and
        usable from a browser and a phone.
\end{itemize}

\section{Failure Cases and Limitations}
\label{sec:disc:limits}

The most common failure modes observed are:

\begin{enumerate}
  \item Visually similar Arabic letters (e.g.\ \textarabic{ك}~vs~\textarabic{ل},
        \textarabic{د}~vs~\textarabic{ذ}, \textarabic{ر}~vs~\textarabic{ز}).
  \item Strong overhead lighting that hides depth information.
  \item Camera distance beyond \SI{1.5}{\metre} where MediaPipe loses
        finger landmarks.
  \item Very low signer counts in \gls{karsl}-502 (3 signers, 8 reps each)
        limit signer-independent generalisation.
  \item MX150 VRAM caps batch size at 32, slowing experimentation.
\end{enumerate}

\section{Comparison with Existing Bilingual SLR}
\label{sec:disc:cmp}

Most prior systems target a single language or only one recognition level.
By covering two languages and two levels in one deployable system, this
work offers a baseline for future bilingual research; the numbers in
\cref{tab:cmp} are competitive despite the small training-time GPU.

\section{Engineering Trade-Offs}
\label{sec:disc:trade}

\begin{table}[H]\centering
  \caption{Key engineering trade-offs in the system.}\label{tab:trade}
  \begin{tabularx}{\textwidth}{l l X}\toprule
    Trade-off                       & Choice           & Justification \\\midrule
    Cloud vs on-device              & Both             & Privacy + reach \\
    Browser MP vs server MP         & Browser          & Privacy + bandwidth \\
    BiLSTM vs Transformer           & BiLSTM           & VRAM, simpler, cuDNN-fast \\
    \texttt{.h5} vs SavedModel      & \texttt{.h5}     & Smaller, Keras-native \\\bottomrule
  \end{tabularx}
\end{table}

\section{Ethical, Social, and Environmental Considerations}
\label{sec:disc:ethics}

The datasets used are gender- and dialect-biased (e.g.\ \gls{karsl}'s
three signers are male Saudi nationals); generalisation to female or
Levantine/Egyptian signers must be validated separately. The on-device
\gls{tflite} option protects user privacy by removing the need to upload
raw video. Approximate energy footprint of training (Kaggle P100 + MX150)
is $\approx XX$~kWh, equivalent to $\approx X$ flights of length
\SI{500}{\kilo\metre}.
